% Annexe E — L'architecture complète.
% Les figures E.1 et E.2 sont générées par notes/gen_carte.py depuis le code ; les tailles sont
% celles du commit de référence, comme tous les chiffres du rapport (chiffres.tex).
\chapter{L'architecture complète}
\label{ann:architecture}

Les sections~\ref{sec:architecture} à~\ref{sec:windows} montrent \demo{} par ses idées : le flux,
le moteur, la décision, les filets. Cette annexe le montre en entier, tel qu'il est écrit. Chaque
boîte de la carte est un fichier réel, chaque nombre est son nombre de lignes, et la carte est
produite par un script qui lit le code (\fichier{notes/gen\_carte.py}) : il échoue si un seul
module n'est rangé nulle part. Elle ne peut donc ni oublier un fichier ni en inventer un.

\begin{table}[htbp]
\centering
\caption{\demo{} en chiffres, au commit de référence \commitref{}}
\label{tab:archi-chiffres}
\small
\begin{tabularx}{\linewidth}{@{}L{6.4cm}Y@{}}
\toprule
\textbf{Mesure} & \textbf{Valeur} \\
\midrule
Modules Python & \nbmodules{}, pour $\approx$ \lignessrc{} lignes \\
Fonctions et méthodes / classes & \nbfonctions{} / \nbclasses{} \\
Liens d'import entre modules du paquet & \nbimports{} \\
Sous-commandes de la ligne de commande & \nbsouscommandes{} \\
Table de règles & \nbregles{} règles, \nbexemples{} exemples-tests, \nboutils{} outils,
  \nbfamilles{} familles \\
Dépendance obligatoire & une seule : \cmd{tomli}, et seulement avant Python 3.11 \\
Dépendances optionnelles & \cmd{winfspy} (montage Windows), \cmd{mitmproxy} (garde de sortie
  HTTP) \\
Tests & \lignestests{} lignes, \fichierstests{} fichiers, \fonctionstests{} fonctions de test \\
\bottomrule
\end{tabularx}
\end{table}

La dépendance unique est un choix : un garde-fou s'installe sur la machine du développeur, à
côté de l'agent, et chaque paquet tiers est une surface de plus à surveiller. Tout le reste ---
l'analyse du shell, la simulation, le hachage du registre, ptrace --- repose sur la bibliothèque
standard de Python. Les deux dépendances optionnelles le sont pour une raison donnée dans
\fichier{pyproject.toml} : si \cmd{winfspy} était obligatoire, une machine sans WinFsp perdrait
\emph{tout} \demo{}, y compris les couches qui n'en ont pas besoin.

\section{La carte}
La figure~\ref{fig:carte} se lit de gauche à droite, dans l'ordre où une commande la traverse :
les portes reçoivent la commande, l'orchestrateur appelle le moteur, le moteur rend un impact, le
décideur tranche, le reçu est écrit. Les flèches portent le nom des appels réels. En bas à
droite, les filets d'exécution, où va une commande supervisée ; à gauche, l'outillage qui
installe, diagnostique et affiche. Les modules en gras sont les pièces centrales de chaque couche.

\begin{sidewaysfigure}
\centering
\resizebox{\textheight}{!}{% Figure E.1 — La carte complète de demo_cli : les 85 modules, groupés par rôle.
% GÉNÉRÉE par notes/gen_carte.py depuis les tailles réelles (lignes par module,
% commit de référence). Ne pas modifier à la main : relancer le script.
\begin{tikzpicture}[x=1cm,y=1cm,
  mod/.style={rectangle,rounded corners=1.2pt,inner xsep=3pt,inner ysep=1.6pt,
              font=\sffamily\tiny,anchor=north west,minimum height=4.2mm},
  liaison/.style={-{Stealth[length=2.4mm]},line width=1.3pt,draw=encre!80},
  appel/.style={font=\ttfamily\tiny,fill=white,draw=encre!35,rounded corners=4pt,
                inner xsep=3pt,inner ysep=1.5pt,text=encre}]
\fill[cPorte!5,rounded corners=4pt,drop shadow={opacity=0.18,shadow xshift=0.7pt,shadow yshift=-0.9pt}] (-0.2,0.8) rectangle (3.35,-4.4);
\fill[cPorte,rounded corners=4pt] (-0.2,0.8) rectangle (3.35,0.14);
\fill[cPorte] (-0.2,0.6000000000000001) rectangle (3.35,0.14);
\draw[cPorte,line width=0.7pt,rounded corners=4pt] (-0.2,0.8) rectangle (3.35,-4.4);
\node[circle,fill=white,text=cPorte,font=\sffamily\tiny\bfseries,inner sep=0pt,minimum size=3.3mm] at (0.09999999999999998,0.5700000000000001) {1};
\node[anchor=west,text=white,font=\sffamily\scriptsize\bfseries,inner sep=0pt] at (0.32,0.5700000000000001) {Portes};
\node[anchor=west,text=white,font=\sffamily\tiny,inner sep=0pt] at (0.32,0.30000000000000004) {7 modules · \num{3549} lignes};
\node[mod,draw=cPorte!85,line width=0.9pt,fill=cPorte!10,text width=3.15cm,path picture={\fill[cPorte!16] ([xshift=1.4pt]path picture bounding box.south west) rectangle ($(path picture bounding box.north west)!0.231!(path picture bounding box.north east)$);\fill[cPorte] (path picture bounding box.south west) rectangle ([xshift=1.4pt]path picture bounding box.north west);}] at (0,0.0) {\bfseries \texttt{hooks/claude\_code.py}\hfill{\color{gris}\num{415}}};
\node[mod,draw=cPorte!45,line width=0.4pt,fill=white,text width=3.15cm,path picture={\fill[cPorte!16] ([xshift=1.4pt]path picture bounding box.south west) rectangle ($(path picture bounding box.north west)!0.297!(path picture bounding box.north east)$);\fill[cPorte] (path picture bounding box.south west) rectangle ([xshift=1.4pt]path picture bounding box.north west);}] at (0,-0.52) {\texttt{hooks/codex.py}\hfill{\color{gris}\num{534}}};
\node[mod,draw=cPorte!45,line width=0.4pt,fill=white,text width=3.15cm,path picture={\fill[cPorte!16] ([xshift=1.4pt]path picture bounding box.south west) rectangle ($(path picture bounding box.north west)!0.104!(path picture bounding box.north east)$);\fill[cPorte] (path picture bounding box.south west) rectangle ([xshift=1.4pt]path picture bounding box.north west);}] at (0,-1.04) {\texttt{hooks/cursor.py}\hfill{\color{gris}\num{187}}};
\node[mod,draw=cPorte!45,line width=0.4pt,fill=white,text width=3.15cm,path picture={\fill[cPorte!16] ([xshift=1.4pt]path picture bounding box.south west) rectangle ($(path picture bounding box.north west)!0.097!(path picture bounding box.north east)$);\fill[cPorte] (path picture bounding box.south west) rectangle ([xshift=1.4pt]path picture bounding box.north west);}] at (0,-1.56) {\texttt{hooks/\_\_init\_\_.py}\hfill{\color{gris}\num{174}}};
\node[mod,draw=cPorte!45,line width=0.4pt,fill=white,text width=3.15cm,path picture={\fill[cPorte!16] ([xshift=1.4pt]path picture bounding box.south west) rectangle ($(path picture bounding box.north west)!0.991!(path picture bounding box.north east)$);\fill[cPorte] (path picture bounding box.south west) rectangle ([xshift=1.4pt]path picture bounding box.north west);}] at (0,-2.08) {\texttt{cli.py}\hfill{\color{gris}\num{1783}}};
\node[mod,draw=cPorte!45,line width=0.4pt,fill=white,text width=3.15cm,path picture={\fill[cPorte!16] ([xshift=1.4pt]path picture bounding box.south west) rectangle ($(path picture bounding box.north west)!0.003!(path picture bounding box.north east)$);\fill[cPorte] (path picture bounding box.south west) rectangle ([xshift=1.4pt]path picture bounding box.north west);}] at (0,-2.6) {\texttt{\_\_main\_\_.py}\hfill{\color{gris}\num{5}}};
\node[mod,draw=cPorte!45,line width=0.4pt,fill=white,text width=3.15cm,path picture={\fill[cPorte!16] ([xshift=1.4pt]path picture bounding box.south west) rectangle ($(path picture bounding box.north west)!0.251!(path picture bounding box.north east)$);\fill[cPorte] (path picture bounding box.south west) rectangle ([xshift=1.4pt]path picture bounding box.north west);}] at (0,-3.12) {\texttt{guarded.py}\hfill{\color{gris}\num{451}}};
\node[font=\sffamily\tiny,text=gris,align=left,anchor=north west,text width=3.2cm] at (0,-3.66) {\textit{piège bash} : \texttt{guard-shell} (dans \texttt{cli.py})};
\fill[cOrch!5,rounded corners=4pt,drop shadow={opacity=0.18,shadow xshift=0.7pt,shadow yshift=-0.9pt}] (4.05,0.8) rectangle (7.6,-2.25);
\fill[cOrch,rounded corners=4pt] (4.05,0.8) rectangle (7.6,0.14);
\fill[cOrch] (4.05,0.6000000000000001) rectangle (7.6,0.14);
\draw[cOrch,line width=0.7pt,rounded corners=4pt] (4.05,0.8) rectangle (7.6,-2.25);
\node[circle,fill=white,text=cOrch,font=\sffamily\tiny\bfseries,inner sep=0pt,minimum size=3.3mm] at (4.35,0.5700000000000001) {2};
\node[anchor=west,text=white,font=\sffamily\scriptsize\bfseries,inner sep=0pt] at (4.57,0.5700000000000001) {Orchestration};
\node[anchor=west,text=white,font=\sffamily\tiny,inner sep=0pt] at (4.57,0.30000000000000004) {4 modules · \num{2476} lignes};
\node[mod,draw=cOrch!85,line width=0.9pt,fill=cOrch!10,text width=3.15cm,path picture={\fill[cOrch!16] ([xshift=1.4pt]path picture bounding box.south west) rectangle ($(path picture bounding box.north west)!0.847!(path picture bounding box.north east)$);\fill[cOrch] (path picture bounding box.south west) rectangle ([xshift=1.4pt]path picture bounding box.north west);}] at (4.25,0.0) {\bfseries \texttt{guard.py}\hfill{\color{gris}\num{1525}}};
\node[mod,draw=cOrch!45,line width=0.4pt,fill=white,text width=3.15cm,path picture={\fill[cOrch!16] ([xshift=1.4pt]path picture bounding box.south west) rectangle ($(path picture bounding box.north west)!0.367!(path picture bounding box.north east)$);\fill[cOrch] (path picture bounding box.south west) rectangle ([xshift=1.4pt]path picture bounding box.north west);}] at (4.25,-0.52) {\texttt{config.py}\hfill{\color{gris}\num{661}}};
\node[mod,draw=cOrch!45,line width=0.4pt,fill=white,text width=3.15cm,path picture={\fill[cOrch!16] ([xshift=1.4pt]path picture bounding box.south west) rectangle ($(path picture bounding box.north west)!0.148!(path picture bounding box.north east)$);\fill[cOrch] (path picture bounding box.south west) rectangle ([xshift=1.4pt]path picture bounding box.north west);}] at (4.25,-1.04) {\texttt{context.py}\hfill{\color{gris}\num{266}}};
\node[mod,draw=cOrch!45,line width=0.4pt,fill=white,text width=3.15cm,path picture={\fill[cOrch!16] ([xshift=1.4pt]path picture bounding box.south west) rectangle ($(path picture bounding box.north west)!0.013!(path picture bounding box.north east)$);\fill[cOrch] (path picture bounding box.south west) rectangle ([xshift=1.4pt]path picture bounding box.north west);}] at (4.25,-1.56) {\texttt{approval.py}\hfill{\color{gris}\num{24}}};
\fill[gris!5,rounded corners=4pt,drop shadow={opacity=0.18,shadow xshift=0.7pt,shadow yshift=-0.9pt}] (4.05,-2.55) rectangle (7.6,-4.4);
\fill[gris,rounded corners=4pt] (4.05,-2.55) rectangle (7.6,-3.21);
\fill[gris] (4.05,-2.75) rectangle (7.6,-3.21);
\draw[gris,line width=0.7pt,rounded corners=4pt] (4.05,-2.55) rectangle (7.6,-4.4);
\node[anchor=west,text=white,font=\sffamily\scriptsize\bfseries,inner sep=0pt] at (4.17,-2.78) {Couche texte (héritée)};
\node[anchor=west,text=white,font=\sffamily\tiny,inner sep=0pt] at (4.17,-3.05) {2 modules · \num{2816} lignes};
\node[mod,draw=gris!45,line width=0.4pt,fill=white,text width=3.15cm,path picture={\fill[gris!16] ([xshift=1.4pt]path picture bounding box.south west) rectangle ($(path picture bounding box.north west)!0.744!(path picture bounding box.north east)$);\fill[gris] (path picture bounding box.south west) rectangle ([xshift=1.4pt]path picture bounding box.north west);}] at (4.25,-3.35) {\texttt{classify.py}\hfill{\color{gris}\num{1339}}};
\node[mod,draw=gris!45,line width=0.4pt,fill=white,text width=3.15cm,path picture={\fill[gris!16] ([xshift=1.4pt]path picture bounding box.south west) rectangle ($(path picture bounding box.north west)!0.821!(path picture bounding box.north east)$);\fill[gris] (path picture bounding box.south west) rectangle ([xshift=1.4pt]path picture bounding box.north west);}] at (4.25,-3.87) {\texttt{targets.py}\hfill{\color{gris}\num{1477}}};
\fill[gris!5,rounded corners=4pt,drop shadow={opacity=0.18,shadow xshift=0.7pt,shadow yshift=-0.9pt}] (-0.2,-4.7) rectangle (7.6,-9.1);
\fill[gris,rounded corners=4pt] (-0.2,-4.7) rectangle (7.6,-5.36);
\fill[gris] (-0.2,-4.9) rectangle (7.6,-5.36);
\draw[gris,line width=0.7pt,rounded corners=4pt] (-0.2,-4.7) rectangle (7.6,-9.1);
\node[circle,fill=white,text=gris,font=\sffamily\tiny\bfseries,inner sep=0pt,minimum size=3.3mm] at (0.09999999999999998,-4.930000000000001) {7};
\node[anchor=west,text=white,font=\sffamily\scriptsize\bfseries,inner sep=0pt] at (0.32,-4.930000000000001) {Installation, diagnostic, affichage};
\node[anchor=west,text=white,font=\sffamily\tiny,inner sep=0pt] at (0.32,-5.2) {8 modules · \num{2651} lignes};
\node[mod,draw=gris!45,line width=0.4pt,fill=white,text width=3.4cm,path picture={\fill[gris!16] ([xshift=1.4pt]path picture bounding box.south west) rectangle ($(path picture bounding box.north west)!0.296!(path picture bounding box.north east)$);\fill[gris] (path picture bounding box.south west) rectangle ([xshift=1.4pt]path picture bounding box.north west);}] at (0,-5.5) {\texttt{doctor.py}\hfill{\color{gris}\num{532}}};
\node[mod,draw=gris!45,line width=0.4pt,fill=white,text width=3.4cm,path picture={\fill[gris!16] ([xshift=1.4pt]path picture bounding box.south west) rectangle ($(path picture bounding box.north west)!0.164!(path picture bounding box.north east)$);\fill[gris] (path picture bounding box.south west) rectangle ([xshift=1.4pt]path picture bounding box.north west);}] at (0,-6.02) {\texttt{deps.py}\hfill{\color{gris}\num{296}}};
\node[mod,draw=gris!45,line width=0.4pt,fill=white,text width=3.4cm,path picture={\fill[gris!16] ([xshift=1.4pt]path picture bounding box.south west) rectangle ($(path picture bounding box.north west)!0.065!(path picture bounding box.north east)$);\fill[gris] (path picture bounding box.south west) rectangle ([xshift=1.4pt]path picture bounding box.north west);}] at (0,-6.54) {\texttt{protection.py}\hfill{\color{gris}\num{117}}};
\node[mod,draw=gris!45,line width=0.4pt,fill=white,text width=3.4cm,path picture={\fill[gris!16] ([xshift=1.4pt]path picture bounding box.south west) rectangle ($(path picture bounding box.north west)!0.349!(path picture bounding box.north east)$);\fill[gris] (path picture bounding box.south west) rectangle ([xshift=1.4pt]path picture bounding box.north west);}] at (0,-7.0600000000000005) {\texttt{lifecycle.py}\hfill{\color{gris}\num{628}}};
\node[mod,draw=gris!45,line width=0.4pt,fill=white,text width=3.4cm,path picture={\fill[gris!16] ([xshift=1.4pt]path picture bounding box.south west) rectangle ($(path picture bounding box.north west)!0.341!(path picture bounding box.north east)$);\fill[gris] (path picture bounding box.south west) rectangle ([xshift=1.4pt]path picture bounding box.north west);}] at (3.95,-5.5) {\texttt{render.py}\hfill{\color{gris}\num{614}}};
\node[mod,draw=gris!45,line width=0.4pt,fill=white,text width=3.4cm,path picture={\fill[gris!16] ([xshift=1.4pt]path picture bounding box.south west) rectangle ($(path picture bounding box.north west)!0.223!(path picture bounding box.north east)$);\fill[gris] (path picture bounding box.south west) rectangle ([xshift=1.4pt]path picture bounding box.north west);}] at (3.95,-6.02) {\texttt{completion.py}\hfill{\color{gris}\num{401}}};
\node[mod,draw=gris!45,line width=0.4pt,fill=white,text width=3.4cm,path picture={\fill[gris!16] ([xshift=1.4pt]path picture bounding box.south west) rectangle ($(path picture bounding box.north west)!0.022!(path picture bounding box.north east)$);\fill[gris] (path picture bounding box.south west) rectangle ([xshift=1.4pt]path picture bounding box.north west);}] at (3.95,-6.54) {\texttt{\_\_init\_\_.py}\hfill{\color{gris}\num{40}}};
\node[mod,draw=gris!45,line width=0.4pt,fill=white,text width=3.4cm,path picture={\fill[gris!16] ([xshift=1.4pt]path picture bounding box.south west) rectangle ($(path picture bounding box.north west)!0.013!(path picture bounding box.north east)$);\fill[gris] (path picture bounding box.south west) rectangle ([xshift=1.4pt]path picture bounding box.north west);}] at (3.95,-7.0600000000000005) {\texttt{version.py}\hfill{\color{gris}\num{23}}};
\node[font=\sffamily\tiny,text=gris,anchor=north west,text width=7.3cm] at (0,-7.65) {\texttt{protection.py} porte le modèle de menace ; \texttt{doctor} et \texttt{deps} disent ce qui manque et la commande qui le corrige ; \texttt{lifecycle} installe et désinstalle.};
\fill[cMoteur!5,rounded corners=4pt,drop shadow={opacity=0.18,shadow xshift=0.7pt,shadow yshift=-0.9pt}] (8.3,0.8) rectangle (16.25,-9.1);
\fill[cMoteur,rounded corners=4pt] (8.3,0.8) rectangle (16.25,0.14);
\fill[cMoteur] (8.3,0.6000000000000001) rectangle (16.25,0.14);
\draw[cMoteur,line width=0.7pt,rounded corners=4pt] (8.3,0.8) rectangle (16.25,-9.1);
\node[circle,fill=white,text=cMoteur,font=\sffamily\tiny\bfseries,inner sep=0pt,minimum size=3.3mm] at (8.600000000000001,0.5700000000000001) {3};
\node[anchor=west,text=white,font=\sffamily\scriptsize\bfseries,inner sep=0pt] at (8.82,0.5700000000000001) {Moteur : comprendre et simuler};
\node[anchor=west,text=white,font=\sffamily\tiny,inner sep=0pt] at (8.82,0.30000000000000004) {46 modules · \num{17019} lignes};
\node[font=\sffamily\tiny\bfseries,text=cMoteur!65!black,anchor=south west,inner sep=0.5pt] at (8.5,-0.26) {le cœur};
\node[mod,draw=cMoteur!85,line width=0.9pt,fill=cMoteur!10,text width=3.6cm,path picture={\fill[cMoteur!16] ([xshift=1.4pt]path picture bounding box.south west) rectangle ($(path picture bounding box.north west)!0.469!(path picture bounding box.north east)$);\fill[cMoteur] (path picture bounding box.south west) rectangle ([xshift=1.4pt]path picture bounding box.north west);}] at (8.5,-0.3) {\bfseries \texttt{model.py}\hfill{\color{gris}\num{845}}};
\node[mod,draw=cMoteur!45,line width=0.4pt,fill=white,text width=3.6cm,path picture={\fill[cMoteur!16] ([xshift=1.4pt]path picture bounding box.south west) rectangle ($(path picture bounding box.north west)!0.086!(path picture bounding box.north east)$);\fill[cMoteur] (path picture bounding box.south west) rectangle ([xshift=1.4pt]path picture bounding box.north west);}] at (8.5,-0.8200000000000001) {\texttt{engine/\_\_init\_\_.py}\hfill{\color{gris}\num{154}}};
\node[mod,draw=cMoteur!45,line width=0.4pt,fill=white,text width=3.6cm,path picture={\fill[cMoteur!16] ([xshift=1.4pt]path picture bounding box.south west) rectangle ($(path picture bounding box.north west)!0.101!(path picture bounding box.north east)$);\fill[cMoteur] (path picture bounding box.south west) rectangle ([xshift=1.4pt]path picture bounding box.north west);}] at (8.5,-1.34) {\texttt{engine/invocation.py}\hfill{\color{gris}\num{182}}};
\node[mod,draw=cMoteur!45,line width=0.4pt,fill=white,text width=3.6cm,path picture={\fill[cMoteur!16] ([xshift=1.4pt]path picture bounding box.south west) rectangle ($(path picture bounding box.north west)!0.119!(path picture bounding box.north east)$);\fill[cMoteur] (path picture bounding box.south west) rectangle ([xshift=1.4pt]path picture bounding box.north west);}] at (8.5,-1.86) {\texttt{engine/words.py}\hfill{\color{gris}\num{214}}};
\node[mod,draw=cMoteur!45,line width=0.4pt,fill=white,text width=3.6cm,path picture={\fill[cMoteur!16] ([xshift=1.4pt]path picture bounding box.south west) rectangle ($(path picture bounding box.north west)!0.169!(path picture bounding box.north east)$);\fill[cMoteur] (path picture bounding box.south west) rectangle ([xshift=1.4pt]path picture bounding box.north west);}] at (8.5,-2.38) {\texttt{engine/outputs.py}\hfill{\color{gris}\num{304}}};
\node[mod,draw=cMoteur!85,line width=0.9pt,fill=cMoteur!10,text width=3.6cm,path picture={\fill[cMoteur!16] ([xshift=1.4pt]path picture bounding box.south west) rectangle ($(path picture bounding box.north west)!0.572!(path picture bounding box.north east)$);\fill[cMoteur] (path picture bounding box.south west) rectangle ([xshift=1.4pt]path picture bounding box.north west);}] at (8.5,-2.9) {\bfseries \texttt{engine/rules.py}\hfill{\color{gris}\num{1030}}};
\node[mod,draw=cMoteur!45,line width=0.4pt,fill=white,text width=3.6cm,path picture={\fill[cMoteur!16] ([xshift=1.4pt]path picture bounding box.south west) rectangle ($(path picture bounding box.north west)!0.422!(path picture bounding box.north east)$);\fill[cMoteur] (path picture bounding box.south west) rectangle ([xshift=1.4pt]path picture bounding box.north west);}] at (8.5,-3.42) {\texttt{engine/grammars.py}\hfill{\color{gris}\num{760}}};
\node[mod,draw=cMoteur!45,line width=0.4pt,fill=white,text width=3.6cm,path picture={\fill[cMoteur!16] ([xshift=1.4pt]path picture bounding box.south west) rectangle ($(path picture bounding box.north west)!0.592!(path picture bounding box.north east)$);\fill[cMoteur] (path picture bounding box.south west) rectangle ([xshift=1.4pt]path picture bounding box.north west);}] at (8.5,-3.94) {\texttt{engine/resolvers.py}\hfill{\color{gris}\num{1065}}};
\node[mod,draw=cMoteur!85,line width=0.9pt,fill=cMoteur!10,text width=3.6cm,path picture={\fill[cMoteur!16] ([xshift=1.4pt]path picture bounding box.south west) rectangle ($(path picture bounding box.north west)!0.631!(path picture bounding box.north east)$);\fill[cMoteur] (path picture bounding box.south west) rectangle ([xshift=1.4pt]path picture bounding box.north west);}] at (8.5,-4.46) {\bfseries \texttt{engine/simulate.py}\hfill{\color{gris}\num{1136}}};
\node[mod,draw=cMoteur!45,line width=0.4pt,fill=white,text width=3.6cm,path picture={\fill[cMoteur!16] ([xshift=1.4pt]path picture bounding box.south west) rectangle ($(path picture bounding box.north west)!0.065!(path picture bounding box.north east)$);\fill[cMoteur] (path picture bounding box.south west) rectangle ([xshift=1.4pt]path picture bounding box.north west);}] at (8.5,-4.9799999999999995) {\texttt{engine/effects.py}\hfill{\color{gris}\num{117}}};
\node[mod,draw=cMoteur!45,line width=0.4pt,fill=white,text width=3.6cm,path picture={\fill[cMoteur!16] ([xshift=1.4pt]path picture bounding box.south west) rectangle ($(path picture bounding box.north west)!0.071!(path picture bounding box.north east)$);\fill[cMoteur] (path picture bounding box.south west) rectangle ([xshift=1.4pt]path picture bounding box.north west);}] at (8.5,-5.5) {\texttt{engine/deferred.py}\hfill{\color{gris}\num{128}}};
\node[font=\sffamily\tiny\bfseries,text=cMoteur!65!black,anchor=south west,inner sep=0.5pt] at (12.350000000000001,-0.26) {12 familles, chargées à la demande};
\node[mod,draw=cMoteur!45,line width=0.4pt,fill=white,text width=3.55cm,path picture={\fill[cMoteur!16] ([xshift=1.4pt]path picture bounding box.south west) rectangle ($(path picture bounding box.north west)!0.788!(path picture bounding box.north east)$);\fill[cMoteur] (path picture bounding box.south west) rectangle ([xshift=1.4pt]path picture bounding box.north west);}] at (12.350000000000001,-0.3) {système {\color{gris}(92 outils)}\hfill{\color{gris}\num{1419}}};
\node[mod,draw=cMoteur!45,line width=0.4pt,fill=white,text width=3.55cm,path picture={\fill[cMoteur!16] ([xshift=1.4pt]path picture bounding box.south west) rectangle ($(path picture bounding box.north west)!0.609!(path picture bounding box.north east)$);\fill[cMoteur] (path picture bounding box.south west) rectangle ([xshift=1.4pt]path picture bounding box.north west);}] at (12.350000000000001,-0.8200000000000001) {envoi {\color{gris}(19 outils)}\hfill{\color{gris}\num{1097}}};
\node[mod,draw=cMoteur!45,line width=0.4pt,fill=white,text width=3.55cm,path picture={\fill[cMoteur!16] ([xshift=1.4pt]path picture bounding box.south west) rectangle ($(path picture bounding box.north west)!0.282!(path picture bounding box.north east)$);\fill[cMoteur] (path picture bounding box.south west) rectangle ([xshift=1.4pt]path picture bounding box.north west);}] at (12.350000000000001,-1.34) {SaaS, migrations {\color{gris}(15 outils)}\hfill{\color{gris}\num{507}}};
\node[mod,draw=cMoteur!45,line width=0.4pt,fill=white,text width=3.55cm,path picture={\fill[cMoteur!16] ([xshift=1.4pt]path picture bounding box.south west) rectangle ($(path picture bounding box.north west)!0.188!(path picture bounding box.north east)$);\fill[cMoteur] (path picture bounding box.south west) rectangle ([xshift=1.4pt]path picture bounding box.north west);}] at (12.350000000000001,-1.86) {PaaS {\color{gris}(8 outils)}\hfill{\color{gris}\num{338}}};
\node[mod,draw=cMoteur!45,line width=0.4pt,fill=white,text width=3.55cm,path picture={\fill[cMoteur!16] ([xshift=1.4pt]path picture bounding box.south west) rectangle ($(path picture bounding box.north west)!0.307!(path picture bounding box.north east)$);\fill[cMoteur] (path picture bounding box.south west) rectangle ([xshift=1.4pt]path picture bounding box.north west);}] at (12.350000000000001,-2.38) {bases de données {\color{gris}(8 outils)}\hfill{\color{gris}\num{552}}};
\node[mod,draw=cMoteur!45,line width=0.4pt,fill=white,text width=3.55cm,path picture={\fill[cMoteur!16] ([xshift=1.4pt]path picture bounding box.south west) rectangle ($(path picture bounding box.north west)!0.514!(path picture bounding box.north east)$);\fill[cMoteur] (path picture bounding box.south west) rectangle ([xshift=1.4pt]path picture bounding box.north west);}] at (12.350000000000001,-2.9) {infrastructure {\color{gris}(6 outils)}\hfill{\color{gris}\num{926}}};
\node[mod,draw=cMoteur!45,line width=0.4pt,fill=white,text width=3.55cm,path picture={\fill[cMoteur!16] ([xshift=1.4pt]path picture bounding box.south west) rectangle ($(path picture bounding box.north west)!0.432!(path picture bounding box.north east)$);\fill[cMoteur] (path picture bounding box.south west) rectangle ([xshift=1.4pt]path picture bounding box.north west);}] at (12.350000000000001,-3.42) {cloud Google, Azure {\color{gris}(5 outils)}\hfill{\color{gris}\num{777}}};
\node[mod,draw=cMoteur!45,line width=0.4pt,fill=white,text width=3.55cm,path picture={\fill[cMoteur!16] ([xshift=1.4pt]path picture bounding box.south west) rectangle ($(path picture bounding box.north west)!0.511!(path picture bounding box.north east)$);\fill[cMoteur] (path picture bounding box.south west) rectangle ([xshift=1.4pt]path picture bounding box.north west);}] at (12.350000000000001,-3.94) {cloud AWS, S3 {\color{gris}(4 outils)}\hfill{\color{gris}\num{920}}};
\node[mod,draw=cMoteur!45,line width=0.4pt,fill=white,text width=3.55cm,path picture={\fill[cMoteur!16] ([xshift=1.4pt]path picture bounding box.south west) rectangle ($(path picture bounding box.north west)!0.784!(path picture bounding box.north east)$);\fill[cMoteur] (path picture bounding box.south west) rectangle ([xshift=1.4pt]path picture bounding box.north west);}] at (12.350000000000001,-4.46) {clients SQL {\color{gris}(4 outils)}\hfill{\color{gris}\num{1412}}};
\node[mod,draw=cMoteur!45,line width=0.4pt,fill=white,text width=3.55cm,path picture={\fill[cMoteur!16] ([xshift=1.4pt]path picture bounding box.south west) rectangle ($(path picture bounding box.north west)!0.164!(path picture bounding box.north east)$);\fill[cMoteur] (path picture bounding box.south west) rectangle ([xshift=1.4pt]path picture bounding box.north west);}] at (12.350000000000001,-4.9799999999999995) {conteneurs {\color{gris}(3 outils)}\hfill{\color{gris}\num{295}}};
\node[mod,draw=cMoteur!45,line width=0.4pt,fill=white,text width=3.55cm,path picture={\fill[cMoteur!16] ([xshift=1.4pt]path picture bounding box.south west) rectangle ($(path picture bounding box.north west)!0.342!(path picture bounding box.north east)$);\fill[cMoteur] (path picture bounding box.south west) rectangle ([xshift=1.4pt]path picture bounding box.north west);}] at (12.350000000000001,-5.5) {git {\color{gris}(1 outil)}\hfill{\color{gris}\num{615}}};
\node[mod,draw=cMoteur!45,line width=0.4pt,fill=white,text width=3.55cm,path picture={\fill[cMoteur!16] ([xshift=1.4pt]path picture bounding box.south west) rectangle ($(path picture bounding box.north west)!0.080!(path picture bounding box.north east)$);\fill[cMoteur] (path picture bounding box.south west) rectangle ([xshift=1.4pt]path picture bounding box.north west);}] at (12.350000000000001,-6.0200000000000005) {Vault {\color{gris}(1 outil)}\hfill{\color{gris}\num{144}}};
\fill[cMoteur!70!black!5,rounded corners=4pt,drop shadow={opacity=0.18,shadow xshift=0.7pt,shadow yshift=-0.9pt}] (8.450000000000001,-6.62) rectangle (16.1,-8.95);
\fill[cMoteur!70!black,rounded corners=4pt] (8.450000000000001,-6.62) rectangle (16.1,-7.08);
\fill[cMoteur!70!black] (8.450000000000001,-6.82) rectangle (16.1,-7.08);
\draw[cMoteur!70!black,line width=0.7pt,rounded corners=4pt,dashed] (8.450000000000001,-6.62) rectangle (16.1,-8.95);
\node[circle,fill=white,text=cMoteur!70!black,font=\sffamily\tiny\bfseries,inner sep=0pt,minimum size=3.3mm,inner xsep=1pt] at (8.750000000000002,-6.8500000000000005) {C3};
\node[anchor=west,text=white,font=\sffamily\scriptsize\bfseries,inner sep=0pt] at (8.97,-6.8500000000000005) {Le frontal PowerShell, sur le même \texttt{World}};
\node[mod,draw=cMoteur!45,line width=0.4pt,fill=white,dashed,text width=2.32cm,path picture={\fill[cMoteur!16] ([xshift=1.4pt]path picture bounding box.south west) rectangle ($(path picture bounding box.north west)!0.381!(path picture bounding box.north east)$);\fill[cMoteur] (path picture bounding box.south west) rectangle ([xshift=1.4pt]path picture bounding box.north west);}] at (8.65,-7.2) {\texttt{ps\_model.py}\hfill{\color{gris}\num{686}}};
\node[mod,draw=cMoteur!45,line width=0.4pt,fill=white,dashed,text width=2.32cm,path picture={\fill[cMoteur!16] ([xshift=1.4pt]path picture bounding box.south west) rectangle ($(path picture bounding box.north west)!0.144!(path picture bounding box.north east)$);\fill[cMoteur] (path picture bounding box.south west) rectangle ([xshift=1.4pt]path picture bounding box.north west);}] at (11.120000000000001,-7.2) {\texttt{ps\_bind.py}\hfill{\color{gris}\num{260}}};
\node[mod,draw=cMoteur!45,line width=0.4pt,fill=white,dashed,text width=2.32cm,path picture={\fill[cMoteur!16] ([xshift=1.4pt]path picture bounding box.south west) rectangle ($(path picture bounding box.north west)!0.322!(path picture bounding box.north east)$);\fill[cMoteur] (path picture bounding box.south west) rectangle ([xshift=1.4pt]path picture bounding box.north west);}] at (13.59,-7.2) {\texttt{ps\_ops.py}\hfill{\color{gris}\num{580}}};
\node[mod,draw=cMoteur!45,line width=0.4pt,fill=white,dashed,text width=3.57cm,path picture={\fill[cMoteur!16] ([xshift=1.4pt]path picture bounding box.south west) rectangle ($(path picture bounding box.north west)!0.235!(path picture bounding box.north east)$);\fill[cMoteur] (path picture bounding box.south west) rectangle ([xshift=1.4pt]path picture bounding box.north west);}] at (8.65,-7.720000000000001) {\texttt{ps\_tables\_data.py}\hfill{\color{gris}\num{423}}};
\node[mod,draw=cMoteur!45,line width=0.4pt,fill=white,dashed,text width=3.57cm,path picture={\fill[cMoteur!16] ([xshift=1.4pt]path picture bounding box.south west) rectangle ($(path picture bounding box.north west)!0.074!(path picture bounding box.north east)$);\fill[cMoteur] (path picture bounding box.south west) rectangle ([xshift=1.4pt]path picture bounding box.north west);}] at (12.370000000000001,-7.720000000000001) {\texttt{rules\_powershell.py}\hfill{\color{gris}\num{133}}};
\node[font=\sffamily\tiny,text=gris,anchor=north west,text width=6.9cm] at (8.65,-8.24) {Tous dans \texttt{engine/} sauf \texttt{ps\_model.py} ; tables générées depuis PowerShell 5.1 : {\cthreecmdlets} cmdlets, {\cthreealias} alias.};
\fill[encre!5,rounded corners=4pt,drop shadow={opacity=0.18,shadow xshift=0.7pt,shadow yshift=-0.9pt}] (16.95,0.8) rectangle (20.25,-1.75);
\fill[encre,rounded corners=4pt] (16.95,0.8) rectangle (20.25,0.14);
\fill[encre] (16.95,0.6000000000000001) rectangle (20.25,0.14);
\draw[encre,line width=0.7pt,rounded corners=4pt] (16.95,0.8) rectangle (20.25,-1.75);
\node[circle,fill=white,text=encre,font=\sffamily\tiny\bfseries,inner sep=0pt,minimum size=3.3mm] at (17.25,0.5700000000000001) {4};
\node[anchor=west,text=white,font=\sffamily\scriptsize\bfseries,inner sep=0pt] at (17.47,0.5700000000000001) {Décider};
\node[anchor=west,text=white,font=\sffamily\tiny,inner sep=0pt] at (17.47,0.30000000000000004) {3 modules · \num{606} lignes};
\node[mod,draw=encre!45,line width=0.4pt,fill=white,text width=3.0cm,path picture={\fill[encre!16] ([xshift=1.4pt]path picture bounding box.south west) rectangle ($(path picture bounding box.north west)!0.063!(path picture bounding box.north east)$);\fill[encre] (path picture bounding box.south west) rectangle ([xshift=1.4pt]path picture bounding box.north west);}] at (17.099999999999998,0.0) {\texttt{engine/verdict.py}\hfill{\color{gris}\num{113}}};
\node[mod,draw=encre!45,line width=0.4pt,fill=white,text width=3.0cm,path picture={\fill[encre!16] ([xshift=1.4pt]path picture bounding box.south west) rectangle ($(path picture bounding box.north west)!0.133!(path picture bounding box.north east)$);\fill[encre] (path picture bounding box.south west) rectangle ([xshift=1.4pt]path picture bounding box.north west);}] at (17.099999999999998,-0.52) {\texttt{policy.py}\hfill{\color{gris}\num{240}}};
\node[mod,draw=encre!85,line width=0.9pt,fill=encre!10,text width=3.0cm,path picture={\fill[encre!16] ([xshift=1.4pt]path picture bounding box.south west) rectangle ($(path picture bounding box.north west)!0.141!(path picture bounding box.north east)$);\fill[encre] (path picture bounding box.south west) rectangle ([xshift=1.4pt]path picture bounding box.north west);}] at (17.099999999999998,-1.04) {\bfseries \texttt{decide.py}\hfill{\color{gris}\num{253}}};
\fill[cRecup!5,rounded corners=4pt,drop shadow={opacity=0.18,shadow xshift=0.7pt,shadow yshift=-0.9pt}] (20.95,0.8) rectangle (24.25,-2.75);
\fill[cRecup,rounded corners=4pt] (20.95,0.8) rectangle (24.25,0.14);
\fill[cRecup] (20.95,0.6000000000000001) rectangle (24.25,0.14);
\draw[cRecup,line width=0.7pt,rounded corners=4pt] (20.95,0.8) rectangle (24.25,-2.75);
\node[circle,fill=white,text=cRecup,font=\sffamily\tiny\bfseries,inner sep=0pt,minimum size=3.3mm] at (21.25,0.5700000000000001) {5};
\node[anchor=west,text=white,font=\sffamily\scriptsize\bfseries,inner sep=0pt] at (21.47,0.5700000000000001) {Récupérer, tracer};
\node[anchor=west,text=white,font=\sffamily\tiny,inner sep=0pt] at (21.47,0.30000000000000004) {5 modules · \num{2598} lignes};
\node[mod,draw=cRecup!45,line width=0.4pt,fill=white,text width=3.0cm,path picture={\fill[cRecup!16] ([xshift=1.4pt]path picture bounding box.south west) rectangle ($(path picture bounding box.north west)!0.633!(path picture bounding box.north east)$);\fill[cRecup] (path picture bounding box.south west) rectangle ([xshift=1.4pt]path picture bounding box.north west);}] at (21.099999999999998,0.0) {\texttt{recovery.py}\hfill{\color{gris}\num{1139}}};
\node[mod,draw=cRecup!45,line width=0.4pt,fill=white,text width=3.0cm,path picture={\fill[cRecup!16] ([xshift=1.4pt]path picture bounding box.south west) rectangle ($(path picture bounding box.north west)!0.118!(path picture bounding box.north east)$);\fill[cRecup] (path picture bounding box.south west) rectangle ([xshift=1.4pt]path picture bounding box.north west);}] at (21.099999999999998,-0.52) {\texttt{checkpoint.py}\hfill{\color{gris}\num{213}}};
\node[mod,draw=cRecup!45,line width=0.4pt,fill=white,text width=3.0cm,path picture={\fill[cRecup!16] ([xshift=1.4pt]path picture bounding box.south west) rectangle ($(path picture bounding box.north west)!0.081!(path picture bounding box.north east)$);\fill[cRecup] (path picture bounding box.south west) rectangle ([xshift=1.4pt]path picture bounding box.north west);}] at (21.099999999999998,-1.04) {\texttt{preview.py}\hfill{\color{gris}\num{146}}};
\node[mod,draw=cRecup!45,line width=0.4pt,fill=white,text width=3.0cm,path picture={\fill[cRecup!16] ([xshift=1.4pt]path picture bounding box.south west) rectangle ($(path picture bounding box.north west)!0.106!(path picture bounding box.north east)$);\fill[cRecup] (path picture bounding box.south west) rectangle ([xshift=1.4pt]path picture bounding box.north west);}] at (21.099999999999998,-1.56) {\texttt{diff.py}\hfill{\color{gris}\num{190}}};
\node[mod,draw=cRecup!85,line width=0.9pt,fill=cRecup!10,text width=3.0cm,path picture={\fill[cRecup!16] ([xshift=1.4pt]path picture bounding box.south west) rectangle ($(path picture bounding box.north west)!0.506!(path picture bounding box.north east)$);\fill[cRecup] (path picture bounding box.south west) rectangle ([xshift=1.4pt]path picture bounding box.north west);}] at (21.099999999999998,-2.08) {\bfseries \texttt{receipts.py}\hfill{\color{gris}\num{910}}};
\fill[cFilet!5,rounded corners=4pt,drop shadow={opacity=0.18,shadow xshift=0.7pt,shadow yshift=-0.9pt}] (16.95,-3.45) rectangle (24.25,-8.2);
\fill[cFilet,rounded corners=4pt] (16.95,-3.45) rectangle (24.25,-4.11);
\fill[cFilet] (16.95,-3.6500000000000004) rectangle (24.25,-4.11);
\draw[cFilet,line width=0.7pt,rounded corners=4pt] (16.95,-3.45) rectangle (24.25,-8.2);
\node[circle,fill=white,text=cFilet,font=\sffamily\tiny\bfseries,inner sep=0pt,minimum size=3.3mm] at (17.25,-3.68) {6};
\node[anchor=west,text=white,font=\sffamily\scriptsize\bfseries,inner sep=0pt] at (17.47,-3.68) {Filets d'exécution};
\node[anchor=west,text=white,font=\sffamily\tiny,inner sep=0pt] at (17.47,-3.95) {10 modules · \num{4822} lignes};
\node[font=\sffamily\tiny\bfseries,text=cFilet!80!black,anchor=south west,inner sep=0.5pt] at (17.099999999999998,-4.38) {Linux (ptrace)};
\node[mod,draw=cFilet!85,line width=0.9pt,fill=cFilet!10,text width=3.0cm,path picture={\fill[cFilet!16] ([xshift=1.4pt]path picture bounding box.south west) rectangle ($(path picture bounding box.north west)!0.347!(path picture bounding box.north east)$);\fill[cFilet] (path picture bounding box.south west) rectangle ([xshift=1.4pt]path picture bounding box.north west);}] at (17.099999999999998,-4.42) {\bfseries \texttt{syscall\_guard.py}\hfill{\color{gris}\num{625}}};
\node[font=\sffamily\tiny\bfseries,text=cFilet!80!black,anchor=south west,inner sep=0.5pt] at (17.099999999999998,-5.32) {réseau (HTTP)};
\node[mod,draw=cFilet!45,line width=0.4pt,fill=white,text width=3.0cm,path picture={\fill[cFilet!16] ([xshift=1.4pt]path picture bounding box.south west) rectangle ($(path picture bounding box.north west)!0.369!(path picture bounding box.north east)$);\fill[cFilet] (path picture bounding box.south west) rectangle ([xshift=1.4pt]path picture bounding box.north west);}] at (17.099999999999998,-5.359999999999999) {\texttt{egress.py}\hfill{\color{gris}\num{665}}};
\node[mod,draw=cFilet!45,line width=0.4pt,fill=white,text width=3.0cm,path picture={\fill[cFilet!16] ([xshift=1.4pt]path picture bounding box.south west) rectangle ($(path picture bounding box.north west)!0.011!(path picture bounding box.north east)$);\fill[cFilet] (path picture bounding box.south west) rectangle ([xshift=1.4pt]path picture bounding box.north west);}] at (17.099999999999998,-5.879999999999999) {\texttt{egress\_addon.py}\hfill{\color{gris}\num{20}}};
\node[font=\sffamily\tiny\bfseries,text=cFilet!80!black,anchor=south west,inner sep=0.5pt] at (21.099999999999998,-4.38) {Windows (WinFsp)};
\node[mod,draw=cFilet!45,line width=0.4pt,fill=white,text width=3.0cm,path picture={\fill[cFilet!16] ([xshift=1.4pt]path picture bounding box.south west) rectangle ($(path picture bounding box.north west)!0.130!(path picture bounding box.north east)$);\fill[cFilet] (path picture bounding box.south west) rectangle ([xshift=1.4pt]path picture bounding box.north west);}] at (21.099999999999998,-4.42) {\texttt{fsguard.py}\hfill{\color{gris}\num{234}}};
\node[mod,draw=cFilet!85,line width=0.9pt,fill=cFilet!10,text width=3.0cm,path picture={\fill[cFilet!16] ([xshift=1.4pt]path picture bounding box.south west) rectangle ($(path picture bounding box.north west)!0.608!(path picture bounding box.north east)$);\fill[cFilet] (path picture bounding box.south west) rectangle ([xshift=1.4pt]path picture bounding box.north west);}] at (21.099999999999998,-4.9399999999999995) {\bfseries \texttt{fsmount.py}\hfill{\color{gris}\num{1095}}};
\node[mod,draw=cFilet!45,line width=0.4pt,fill=white,text width=3.0cm,path picture={\fill[cFilet!16] ([xshift=1.4pt]path picture bounding box.south west) rectangle ($(path picture bounding box.north west)!0.263!(path picture bounding box.north east)$);\fill[cFilet] (path picture bounding box.south west) rectangle ([xshift=1.4pt]path picture bounding box.north west);}] at (21.099999999999998,-5.46) {\texttt{fspassthrough.py}\hfill{\color{gris}\num{473}}};
\node[mod,draw=cFilet!45,line width=0.4pt,fill=white,text width=3.0cm,path picture={\fill[cFilet!16] ([xshift=1.4pt]path picture bounding box.south west) rectangle ($(path picture bounding box.north west)!0.098!(path picture bounding box.north east)$);\fill[cFilet] (path picture bounding box.south west) rectangle ([xshift=1.4pt]path picture bounding box.north west);}] at (21.099999999999998,-5.98) {\texttt{fsreport.py}\hfill{\color{gris}\num{177}}};
\node[mod,draw=cFilet!45,line width=0.4pt,fill=white,text width=3.0cm,path picture={\fill[cFilet!16] ([xshift=1.4pt]path picture bounding box.south west) rectangle ($(path picture bounding box.north west)!0.096!(path picture bounding box.north east)$);\fill[cFilet] (path picture bounding box.south west) rectangle ([xshift=1.4pt]path picture bounding box.north west);}] at (21.099999999999998,-6.5) {\texttt{mountstate.py}\hfill{\color{gris}\num{172}}};
\node[mod,draw=cFilet!45,line width=0.4pt,fill=white,text width=3.0cm,path picture={\fill[cFilet!16] ([xshift=1.4pt]path picture bounding box.south west) rectangle ($(path picture bounding box.north west)!0.529!(path picture bounding box.north east)$);\fill[cFilet] (path picture bounding box.south west) rectangle ([xshift=1.4pt]path picture bounding box.north west);}] at (21.099999999999998,-7.02) {\texttt{protect.py}\hfill{\color{gris}\num{953}}};
\node[mod,draw=cFilet!45,line width=0.4pt,fill=white,text width=3.0cm,path picture={\fill[cFilet!16] ([xshift=1.4pt]path picture bounding box.south west) rectangle ($(path picture bounding box.north west)!0.227!(path picture bounding box.north east)$);\fill[cFilet] (path picture bounding box.south west) rectangle ([xshift=1.4pt]path picture bounding box.north west);}] at (21.099999999999998,-7.54) {\texttt{schedule.py}\hfill{\color{gris}\num{408}}};
\draw[trait,rounded corners=3pt,fill=white] (16.95,-8.5) rectangle (24.25,-9.1);
\node[mod,draw=encre!45,line width=0.4pt,fill=white,text width=1.15cm,path picture={\fill[encre!16] ([xshift=1.4pt]path picture bounding box.south west) rectangle ($(path picture bounding box.north west)!0.500!(path picture bounding box.north east)$);\fill[encre] (path picture bounding box.south west) rectangle ([xshift=1.4pt]path picture bounding box.north west);}] at (17.099999999999998,-8.59) {\texttt{module}\hfill{\color{gris}\num{900}}};
\node[font=\sffamily\tiny,text=encre,anchor=west] at (18.499999999999996,-8.8) {barre : taille relative};
\node[mod,draw=encre!85,line width=0.9pt,fill=encre!10,text width=0.75cm,path picture={\fill[encre] (path picture bounding box.south west) rectangle ([xshift=1.4pt]path picture bounding box.north west);}] at (20.749999999999996,-8.59) {\bfseries \texttt{clé}};
\node[font=\sffamily\tiny,text=encre,anchor=west] at (21.65,-8.8) {pièce centrale};
\node[mod,draw=cMoteur!45,line width=0.4pt,fill=white,dashed,text width=0.45cm,path picture={\fill[cMoteur] (path picture bounding box.south west) rectangle ([xshift=1.4pt]path picture bounding box.north west);}] at (23.299999999999997,-8.59) {\texttt{C3}};
\draw[liaison] (3.39,0.47000000000000003) -- (4.01,0.47000000000000003);
\node[appel] at (3.7,1.1) {evaluate};
\draw[encre!35,line width=0.4pt] (3.7,0.9700000000000001) -- (3.7,0.53);
\draw[liaison] (7.64,0.47000000000000003) -- (8.260000000000002,0.47000000000000003);
\node[appel] at (7.95,1.1) {\_engine\_predict};
\draw[encre!35,line width=0.4pt] (7.95,0.9700000000000001) -- (7.95,0.53);
\draw[liaison] (16.29,0.47000000000000003) -- (16.91,0.47000000000000003);
\node[appel] at (16.6,1.1) {Impact};
\draw[encre!35,line width=0.4pt] (16.6,0.9700000000000001) -- (16.6,0.53);
\draw[liaison] (20.29,0.47000000000000003) -- (20.91,0.47000000000000003);
\node[appel] at (20.6,1.1) {\_record};
\draw[encre!35,line width=0.4pt] (20.6,0.9700000000000001) -- (20.6,0.53);
\draw[liaison] (19.95,-1.75) -- (19.95,-3.45);
\node[appel,anchor=east] at (19.8,-2.6) {SUPERVISED : sh -g};
\draw[liaison] (22.849999999999998,-3.45) -- (22.849999999999998,-2.75);
\node[appel,anchor=west,align=left] at (23.0,-3.1) {copie avant\\destruction};
\end{tikzpicture}
}
\caption[La carte complète de \demo{} : les \nbmodules{} modules, groupés par rôle]{La carte complète de \demo{} : les \nbmodules{} modules, groupés par rôle, avec leur
nombre de lignes au commit \commitref{}. Les familles du moteur sont regroupées (un fichier de
règles, un fichier d'opérations, parfois un fichier de données générées) ; le détail est au
tableau~\ref{tab:familles}. C3, en tirets, est décrit à la section~\ref{sec:c3}.}
\label{fig:carte}
\end{sidewaysfigure}

\section{Où sont les lignes}
La figure~\ref{fig:couches} répond à une question simple : où est passé l'effort ? Le moteur porte
43\,\% du code, et les trois cinquièmes du moteur sont dans ses familles ---
la connaissance de ce que font \nboutils{} outils, écrite comme des données et vérifiée par
leurs exemples. Le décideur, lui, tient en quelques centaines de lignes : c'est voulu. Une
décision doit pouvoir se relire en entier ; ce qui est long, c'est de \emph{savoir} ce qu'une
commande fait, pas de décider quoi en faire.

\begin{figure}[htbp]
\centering
% Figure E.2 — Où sont les lignes. GÉNÉRÉE par notes/gen_carte.py.
\begin{tikzpicture}[x=1cm,y=1cm]
\node[anchor=east,font=\sffamily\small] at (-0.15,0.0) {Portes};
\fill[encre!55,rounded corners=1pt] (0,-0.2) rectangle (2.878,0.2);
\node[anchor=west,font=\sffamily\footnotesize] at (2.998,0.0) {\num{3549} \textcolor{gris}{· 7 modules · 10\,\%}};
\node[anchor=east,font=\sffamily\small] at (-0.15,-0.62) {Orchestration};
\fill[encre!55,rounded corners=1pt] (0,-0.8200000000000001) rectangle (2.008,-0.42);
\node[anchor=west,font=\sffamily\footnotesize] at (2.128,-0.62) {\num{2476} \textcolor{gris}{· 4 modules · 7\,\%}};
\node[anchor=east,font=\sffamily\small] at (-0.15,-1.24) {Couche texte (héritée)};
\fill[encre!55,rounded corners=1pt] (0,-1.44) rectangle (2.284,-1.04);
\node[anchor=west,font=\sffamily\footnotesize] at (2.404,-1.24) {\num{2816} \textcolor{gris}{· 2 modules · 8\,\%}};
\node[anchor=east,font=\sffamily\small] at (-0.15,-1.8599999999999999) {Moteur : cœur};
\fill[cMoteur,rounded corners=1pt] (0,-2.06) rectangle (4.813,-1.66);
\node[anchor=west,font=\sffamily\footnotesize] at (4.933,-1.8599999999999999) {\num{5935} \textcolor{gris}{· 11 modules · 16\,\%}};
\node[anchor=east,font=\sffamily\small] at (-0.15,-2.48) {Moteur : 12 familles};
\fill[cMoteur,rounded corners=1pt] (0,-2.68) rectangle (7.300,-2.28);
\node[anchor=west,font=\sffamily\footnotesize] at (7.420,-2.48) {\num{9002} \textcolor{gris}{· 30 modules · 25\,\%}};
\node[anchor=east,font=\sffamily\small] at (-0.15,-3.1) {Moteur : frontal PowerShell};
\fill[cMoteur,rounded corners=1pt] (0,-3.3000000000000003) rectangle (1.688,-2.9);
\node[anchor=west,font=\sffamily\footnotesize] at (1.808,-3.1) {\num{2082} \textcolor{gris}{· 5 modules · 6\,\%}};
\node[anchor=east,font=\sffamily\small] at (-0.15,-3.7199999999999998) {Décider};
\fill[encre!55,rounded corners=1pt] (0,-3.92) rectangle (0.491,-3.5199999999999996);
\node[anchor=west,font=\sffamily\footnotesize] at (0.611,-3.7199999999999998) {\num{606} \textcolor{gris}{· 3 modules · 2\,\%}};
\node[anchor=east,font=\sffamily\small] at (-0.15,-4.34) {Récupérer, tracer};
\fill[encre!55,rounded corners=1pt] (0,-4.54) rectangle (2.107,-4.14);
\node[anchor=west,font=\sffamily\footnotesize] at (2.227,-4.34) {\num{2598} \textcolor{gris}{· 5 modules · 7\,\%}};
\node[anchor=east,font=\sffamily\small] at (-0.15,-4.96) {Filets d'exécution};
\fill[encre!55,rounded corners=1pt] (0,-5.16) rectangle (3.910,-4.76);
\node[anchor=west,font=\sffamily\footnotesize] at (4.030,-4.96) {\num{4822} \textcolor{gris}{· 10 modules · 13\,\%}};
\node[anchor=east,font=\sffamily\small] at (-0.15,-5.58) {Installation, diagnostic};
\fill[encre!55,rounded corners=1pt] (0,-5.78) rectangle (2.150,-5.38);
\node[anchor=west,font=\sffamily\footnotesize] at (2.270,-5.58) {\num{2651} \textcolor{gris}{· 8 modules · 7\,\%}};
\end{tikzpicture}

\caption{Lignes de code par couche, au commit \commitref{} (même découpage que la
figure~\ref{fig:carte})}
\label{fig:couches}
\end{figure}

\begin{table}[htbp]
\centering
\caption{Les \nbfamilles{} familles du moteur, chargées à la demande}
\label{tab:familles}
\small
\begin{tabularx}{\linewidth}{@{}L{3.3cm}Yrr@{}}
\toprule
\textbf{Famille} & \textbf{Exemples d'outils} & \textbf{Outils} & \textbf{Lignes} \\
\midrule
Système & \cmd{systemctl}, \cmd{crontab}, \cmd{chmod}, \cmd{useradd}, \cmd{reg}, \cmd{schtasks} & 92 & \num{1419} \\
Clients SQL & \cmd{sqlite3}, \cmd{psql}, \cmd{mysql}, \cmd{mariadb} & 4 & \num{1412} \\
Envoi & \cmd{curl}, \cmd{gh}, \cmd{ssh}, \cmd{scp}, \cmd{npm}, \cmd{cargo}, \cmd{sendmail} & 19 & \num{1097} \\
Infrastructure & \cmd{terraform}, \cmd{tofu}, \cmd{kubectl}, \cmd{helm}, \cmd{pulumi}, \cmd{cdk} & 6 & 926 \\
Cloud AWS, S3 & \cmd{aws}, \cmd{s3cmd}, \cmd{rclone}, \cmd{mc} & 4 & 920 \\
Cloud Google, Azure & \cmd{gcloud}, \cmd{gsutil}, \cmd{az}, \cmd{azcopy}, \cmd{railway} & 5 & 777 \\
git & \cmd{git} & 1 & 615 \\
Bases de données & \cmd{redis-cli}, \cmd{mongosh}, \cmd{dropdb}, \cmd{pg\_restore}, \cmd{sqlcmd} & 8 & 552 \\
SaaS, migrations & \cmd{shopify}, \cmd{wp}, \cmd{stripe}, \cmd{alembic}, \cmd{prisma}, \cmd{flyway} & 15 & 507 \\
PaaS & \cmd{heroku}, \cmd{fly}, \cmd{vercel}, \cmd{netlify}, \cmd{wrangler}, \cmd{firebase} & 8 & 338 \\
Conteneurs & \cmd{docker}, \cmd{podman}, \cmd{docker-compose} & 3 & 295 \\
Vault & \cmd{vault} & 1 & 144 \\
\midrule
\multicolumn{2}{@{}l}{Règles du cœur (\cmd{rm}, \cmd{mv}, \cmd{cp}, \cmd{sed}, \cmd{tar}, \cmd{find}, \cmd{rsync}\ldots)} & 48 & --- \\
\bottomrule
\end{tabularx}
\end{table}

Une famille n'est importée que lorsqu'une ligne nomme l'un de ses outils
(\cmd{LAZY\_FAMILIES} dans \fichier{engine/rules.py}) : un appel du \textit{hook} qui ne
rencontre que \cmd{ls} et \cmd{rm} ne paie ni les grammaires de \cmd{aws}, ni les
\num{1419} lignes de la famille système. C'est ce qui permet au moteur de grossir famille par
famille sans sortir de son budget de temps. Les tables les plus longues (sous-commandes de
\cmd{gh}, de \cmd{npm}, de \cmd{stripe}, de \cmd{vault}\ldots) ne sont pas écrites à la main :
elles sont générées depuis les fichiers des outils eux-mêmes, et le fichier généré le dit.

\section{Le trajet d'une commande dans le code}
La figure~\ref{fig:flux} du chapitre~3 montre le flux comme une idée. La
figure~\ref{fig:trajet} montre le même flux comme du code : les fonctions réelles, dans l'ordre
où \cmd{Guard.evaluate} les appelle, pour une commande Bash reçue par le \textit{hook} de
Claude Code. Les autres portes (Codex, Cursor, le piège bash, \cmd{demo\_cli check}) entrent à
l'étape~2 par la même fonction : il n'y a qu'un chemin de décision.

\begin{figure}[htbp]
\centering
% Figure E.3 — Le trajet d'une commande dans le code : les fonctions réelles, dans l'ordre
% où Guard.evaluate les appelle (guard.py, commit de référence).
\begin{tikzpicture}[x=1cm,y=1cm,
  etape/.style={circle,fill=encre,text=white,font=\sffamily\scriptsize\bfseries,
                inner sep=0pt,minimum size=4.6mm},
  modu/.style={rounded corners=1.5pt,font=\ttfamily\scriptsize,text=white,inner sep=2.5pt,
               minimum height=4.6mm,minimum width=33mm,anchor=west},
  fonc/.style={font=\ttfamily\scriptsize,anchor=west,text=encre},
  quoi/.style={font=\sffamily\scriptsize,anchor=west,text=gris,text width=5.95cm}]
  \def\pas{0.82}
  \draw[trait,line width=1.4pt] (0,0) -- (0,-11*\pas);
  \foreach \i/\m/\c/\f/\q in {
    1/{hooks/claude\_code.py}/cPorte/{run\_pretooluse}/{lit l'appel d'outil de l'agent (JSON sur l'entrée standard)},
    2/{guard.py}/cOrch/{Guard.evaluate}/{point d'entrée unique, pour toutes les portes},
    3/{guard.py}/cOrch/{\_early\_refusal}/{configuration illisible ; sous-commande d'état lancée par l'agent},
    4/{classify.py}/gris/{classify\_pipeline}/{la lecture historique du texte, gardée pour le chemin hérité},
    5/{engine/}/cMoteur/{predict $\rightarrow$ simulate $\rightarrow$ verdict}/{analyse, invocations, règles, FS virtuel : l'impact et son jugement},
    6/{guard.py}/cOrch/{\_control\_plane\_refusal}/{un fichier de \texttt{demo\_cli} touché : ESCALATE, sans discussion},
    7/{guard.py}/cOrch/{\_engine\_facts | \_unread\_facts}/{chemin moteur, ou ligne non lue : filet ou point de contrôle},
    8/{recovery.py}/cRecup/{snapshot\_set}/{la copie, tout ou rien, \emph{avant} la décision},
    9/{context.py}/cOrch/{build\_context}/{environnement déclaré d'abord ; comparaison avec l'intention},
    10/{decide.py}/encre/{decide}/{le décideur unique, pur, avec \texttt{policy.py}},
    11/{receipts.py}/cRecup/{append\_receipt}/{le reçu chaîné (SHA-256), pour \emph{chaque} décision},
    12/{hooks/claude\_code.py}/cPorte/{\_emit}/{allow, ask, deny --- ou réécriture vers \texttt{demo\_cli sh -g}}}
  {
    \node[etape] at (0,{-(\i-1)*\pas}) {\i};
    \node[modu,fill=\c] at (0.45,{-(\i-1)*\pas}) {\m};
    \node[fonc] at (4.05,{-(\i-1)*\pas}) {\f};
    \node[quoi] at (9.45,{-(\i-1)*\pas}) {\q};
  }
\end{tikzpicture}

\caption{Le trajet d'une commande dans le code, du \textit{hook} au reçu (fonctions réelles de
\fichier{guard.py} et des modules qu'il appelle)}
\label{fig:trajet}
\end{figure}

\section{Les dépendances vont dans un seul sens}
La carte a été vérifiée sur les \nbimports{} liens d'import réels du paquet. Trois propriétés
en ressortent, et elles ne doivent rien au hasard :
\begin{itemize}
  \item \textbf{Le décideur n'importe presque rien.} \fichier{decide.py} importe
    \fichier{policy.py} et le type de classification, et rien d'autre : aucun accès au disque,
    au réseau ou au registre. C'est ce qui le rend pur, testable en isolation, et incapable de
    « mentir » sur une capture : il ne peut que lire le fait qu'on lui donne.
  \item \textbf{Le moteur ne connaît pas l'orchestrateur.} Aucun module de \fichier{engine/}
    n'importe \fichier{guard.py} ni une porte. Le moteur peut donc être exercé seul : les
    \nbexemples{} exemples de la table le traversent directement dans les tests, et le
    préfiltre du piège shell est généré depuis sa table de règles.
  \item \textbf{La politique est à un seul endroit.} \fichier{decide.py},
    \fichier{engine/verdict.py}, \fichier{guard.py} et le \textit{hook} de Claude Code lisent
    \fichier{policy.py} pour toute question sur une surface distante. Changer la décision d'une
    surface, c'est changer une ligne, une fois.
\end{itemize}

\section{Carte des modules}
Le tableau~\ref{tab:modules} donne le rôle de chaque module hors familles (les familles sont au
tableau~\ref{tab:familles}), dans l'ordre de la carte.

\begingroup
\small\linespread{1.0}\selectfont
\setlength{\LTpre}{6pt}\setlength{\LTpost}{6pt}
\begin{longtable}{@{}L{4.1cm}r@{\hspace{8pt}}L{9.85cm}@{}}
\caption{Carte des modules : rôle et taille (lignes, commit \commitref{})}
\label{tab:modules}\\
\toprule
\textbf{Module} & \textbf{Lignes} & \textbf{Rôle} \\
\midrule
\endfirsthead
\toprule
\textbf{Module} & \textbf{Lignes} & \textbf{Rôle} \\
\midrule
\endhead
\midrule
\multicolumn{3}{r}{\footnotesize\itshape suite page suivante}\\
\endfoot
\bottomrule
\endlastfoot
\multicolumn{3}{@{}l}{\textcolor{cPorte}{\sffamily\bfseries 1 --- Portes}} \\
\fichier{hooks/claude\_code.py} & 415 & \textit{Hook} PreToolUse de Claude Code : lit l'appel
  d'outil, appelle le garde, répond allow, ask ou deny, ou réécrit la commande vers
  \cmd{demo\_cli sh -g}. \\
\fichier{hooks/codex.py} & 524 & \textit{Hook} de Codex : la même décision, exprimée par refus
  seul. \\
\fichier{hooks/cursor.py} & 187 & \textit{Hook} \cmd{beforeShellExecution} de Cursor, fermé en
  cas d'échec. \\
\fichier{hooks/\_\_init\_\_.py} & 174 & Le socle commun des intégrations. \\
\fichier{cli.py} & \num{1783} & La ligne de commande : \nbsouscommandes{} sous-commandes
  (\cmd{check}, \cmd{undo}, \cmd{verify}, \cmd{sh}, \cmd{setup}\ldots) et le piège bash
  (\cmd{guard-shell}). Aucune logique de sécurité : elle analyse, appelle, affiche. \\
\fichier{\_\_main\_\_.py} & 5 & \cmd{python -m demo\_cli}. \\
\fichier{guarded.py} & 451 & \cmd{demo\_cli guarded} : lance l'agent avec toutes les couches
  qui peuvent l'être, et dit lesquelles manquent. \\
\multicolumn{3}{@{}l}{\textcolor{cOrch}{\sffamily\bfseries 2 --- Orchestration}} \\
\fichier{guard.py} & \num{1461} & L'orchestrateur : \cmd{Guard.evaluate} enchaîne les étapes
  nommées (figure~\ref{fig:trajet}) ; \cmd{evaluate\_file\_edit} pour la porte fichier. \\
\fichier{config.py} & 661 & \fichier{.demo\_cli.toml} : mode (observation ou application),
  cibles déclarées, espace de travail propre au projet. \\
\fichier{context.py} & 266 & Le contexte : empreinte, environnement (déclaré d'abord, deviné en
  dernier), comparaison avec l'intention déclarée. \\
\fichier{approval.py} & 24 & L'autorisation structurelle : un jeton signé par une clé que
  l'agent ne détient pas. \\
\multicolumn{3}{@{}l}{\textcolor{gris}{\sffamily\bfseries Couche texte (héritée)}} \\
\fichier{classify.py} & \num{1339} & La lecture historique du texte d'une commande, gardée pour
  le chemin hérité ; elle porte aussi les constantes de dialecte. \\
\fichier{targets.py} & \num{1477} & La cible réelle d'une commande (fichier, SQLite,
  PostgreSQL), sans jamais de cible par défaut. \\
\multicolumn{3}{@{}l}{\textcolor{cMoteur!80!black}{\sffamily\bfseries 3 --- Moteur : le cœur}} \\
\fichier{model.py} & 845 & L'analyse POSIX unique : mots et citations, substitutions
  imbriquées, \textit{heredocs}, enveloppes. \\
\fichier{engine/\_\_init\_\_.py} & 98 & Le point d'entrée : \cmd{predict}, son budget de temps,
  la prédiction écrite dans le reçu. \\
\fichier{engine/invocation.py} & 182 & De l'analyse aux invocations : une par commande qui
  s'exécute, avec ses enveloppes et ses scripts imbriqués. \\
\fichier{engine/words.py} & 214 & L'expansion des mots comme bash, sur les mots ; l'inconnu est
  levé, jamais deviné. \\
\fichier{engine/outputs.py} & 304 & Ce qu'impriment quelques lectures, pour lire
  \cmd{rm \$(cat liste.txt)}. \\
\fichier{engine/rules.py} & \num{1030} & La table de règles comme données ; le chargement des
  familles ; le préfiltre du piège shell. \\
\fichier{engine/grammars.py} & 760 & Les grammaires d'options générées depuis l'aide des outils.
  \\
\fichier{engine/resolvers.py} & \num{1065} & Les outils qu'un gabarit ne peut pas exprimer
  (\cmd{cp}, \cmd{sed}, \cmd{tar}, \cmd{rsync}, \cmd{find}, \cmd{dd}\ldots). \\
\fichier{engine/simulate.py} & \num{1110} & La simulation sur un système de fichiers virtuel :
  ensembles d'états, origine du contenu, budget. \\
\fichier{engine/effects.py} & 111 & Le vocabulaire des effets
  (tableau~\ref{tab:effets}). \\
\fichier{engine/deferred.py} & 128 & Les lieux où du code s'exécute plus tard sans être nommé
  (profil du shell, \textit{hooks} git, tâches planifiées). \\
\multicolumn{3}{@{}l}{\textcolor{encre}{\sffamily\bfseries 4 --- Décider}} \\
\fichier{engine/verdict.py} & 113 & Le jugement d'un impact, dans l'ordre convenu. \\
\fichier{policy.py} & 238 & La politique unique : décision par surface, canal ask ou deny,
  indices de récupération. \\
\fichier{decide.py} & 251 & Le décideur unique, pur, à l'ordre de branches fixe. \\
\multicolumn{3}{@{}l}{\textcolor{cRecup}{\sffamily\bfseries 5 --- Récupérer, tracer}} \\
\fichier{recovery.py} & \num{1139} & Les captures tout-ou-rien, les groupes, la restauration
  (\cmd{undo}). \\
\fichier{checkpoint.py} & 213 & Le point de contrôle de tout le projet, optionnel et borné. \\
\fichier{preview.py} & 146 & L'aperçu SQL (DRY\_RUN) : lignes touchées et échantillon. \\
\fichier{diff.py} & 190 & Ce qui a changé depuis le dernier point de récupération. \\
\fichier{receipts.py} & 903 & Le registre chaîné : ajout sous verrou, vérification de la
  chaîne, liens entre chaînes. \\
\multicolumn{3}{@{}l}{\textcolor{cFilet}{\sffamily\bfseries 6 --- Filets d'exécution}} \\
\fichier{syscall\_guard.py} & 528 & Linux : le garde ptrace (section~\ref{sec:filets}). \\
\fichier{fsguard.py} & 234 & Windows : le jugement --- quelles opérations de fichier
  détruisent. \\
\fichier{fsmount.py} & \num{1095} & Windows : le système de fichiers WinFsp et son montage. \\
\fichier{fspassthrough.py} & 473 & Windows : le passe-plat sur les fichiers réels. \\
\fichier{fsreport.py} & 177 & Dire à l'agent ce que le montage a capturé. \\
\fichier{mountstate.py} & 172 & Le montage tourne-t-il \emph{vraiment} ? (« installé mais
  inerte ») \\
\fichier{protect.py} & 953 & Déplacer le projet vers \cmd{.real}, le verrouiller --- et
  l'inverse. \\
\fichier{schedule.py} & 408 & La tâche d'ouverture de session durcie
  (section~\ref{sec:wsec}). \\
\fichier{egress.py} & 665 & Le garde de sortie : les écritures vers des API SaaS, vues au
  niveau du réseau. \\
\fichier{egress\_addon.py} & 20 & Le chargeur de ce garde pour \cmd{mitmproxy}. \\
\multicolumn{3}{@{}l}{\textcolor{gris}{\sffamily\bfseries 7 --- Installation, diagnostic, affichage}} \\
\fichier{doctor.py} & 532 & Le diagnostic : \textit{hooks}, montage, tâche, dépendances. \\
\fichier{deps.py} & 296 & Ce qui manque, et la commande exacte qui le corrige. \\
\fichier{protection.py} & 115 & Le modèle de menace, et l'état des protections actives. \\
\fichier{lifecycle.py} & 628 & \cmd{setup} et \cmd{teardown}, étape par étape. \\
\fichier{render.py} & 611 & L'affichage : couleurs des décisions, \cmd{NO\_COLOR}. \\
\fichier{completion.py} & 401 & La complétion pour bash, zsh, fish et PowerShell. \\
\fichier{\_\_init\_\_.py} & 40 & L'API publique et l'invariant. \\
\fichier{version.py} & 23 & La version, à un seul endroit. \\
\end{longtable}
\endgroup

\begin{pourquoi}[Ce que la carte raconte]
\demo{} est parti d'une bêta de \betalignessrc{} lignes : une classification par expressions
régulières et cinq décisions. La carte montre ce qu'il a fallu construire autour pour tenir
une seule promesse --- « une action qui modifie doit être récupérable et correspondre au
contexte déclaré » : un analyseur, un simulateur, \nboutils{} outils décrits comme des données,
un décideur unique, deux filets d'exécution, un registre dont toute retouche se voit, et l'outillage qui
installe tout cela proprement sur deux systèmes. Chaque boîte répond à une question précise ;
aucune n'est là pour faire nombre.
\end{pourquoi}
